%==============================================================================
% ModelCrew · CUMCM 国赛论文模板（中文，ctex，自包含）
%   编译：  latexmk -xelatex 6_paper.tex   （含中文，建议 XeLaTeX）
%   若本机装有官方 cumcmthesis.cls，可把 \documentclass 换成：
%       \documentclass{cumcmthesis}
%   占位符约定见 templates/README.md —— Writer 只替换 %%...%% 内容，不动排版骨架。
%   注意：国赛真实提交需"承诺书 + 编号页"，下方 %%COMMITMENT%% 为占位，按当年官方模板替换。
%==============================================================================
\documentclass[12pt,a4paper]{ctexart}

\usepackage[margin=2.5cm]{geometry}
\usepackage{amsmath,amssymb}
\usepackage{graphicx}
\usepackage{booktabs}
\usepackage{titlesec}
\usepackage{abstract}
\usepackage[hidelinks]{hyperref}

\title{\textbf{%%TITLE%%}}
\author{}
\date{}

\begin{document}

%------------------------------------------------------------------ 承诺页（占位）
% 国赛要求的"诚信承诺书"，按当年官方模板替换本段
\thispagestyle{empty}
\noindent %%COMMITMENT%%
\newpage

%------------------------------------------------------------------ 摘要页
\maketitle
\vspace{-2em}
\begin{center}\textbf{\large 摘\quad 要}\end{center}
\noindent
%%SUMMARY%%

\vspace{1em}
\noindent\textbf{关键词：} %%KEYWORDS%%

\newpage
\tableofcontents
\newpage

%------------------------------------------------------------------ 正文
\section{问题重述}
%%SEC_RESTATE%%

\section{问题分析与基本假设}
%%SEC_ASSUMPTIONS%%

\section{符号说明}
%%SEC_NOTATION%%

\section{模型的建立}
%%SEC_MODEL%%

\section{模型的求解与结果}
%%SEC_SOLUTION%%

\section{灵敏度分析}
%%SEC_SENSITIVITY%%

\section{模型的评价}
% 优点 / 局限（如实写入 Critic 审计的 ⚠️/❌）/ 改进方向
%%SEC_EVALUATION%%

\section{结论}
%%SEC_CONCLUSION%%

\begin{thebibliography}{9}
%%REFERENCES%%
\end{thebibliography}

\appendix
\section{附录}
%%APPENDIX%%

\end{document}
